\documentclass[a4paper]{arrowhead}

\usepackage[yyyymmdd]{datetime}
\usepackage{etoolbox}
\usepackage[utf8]{inputenc}
\usepackage[T1]{fontenc}
\usepackage{multirow}
\usepackage{longtable}
\usepackage{tikz}
\usetikzlibrary{positioning,shapes.multipart,arrows.meta}

\renewcommand{\dateseparator}{-}
\setlength{\emergencystretch}{3em}
\lstset{language=HTTP,breakatwhitespace=false,columns=fullflexible}

%% Special references
\newcommand{\mref}[2][]{{\textcolor{ArrowheadPurple}{\hyperref[sec:model:#2]{\ifstrempty{#1}{#2}{#1}}}}}
\newcommand{\pref}[1]{{\textcolor{ArrowheadGrey}{\hyperref[sec:model:primitives:#1]{#1}}}}
%%

\begin{document}
\sloppy

%% Arrowhead Document Properties
\ArrowheadTitle{GraphDBManager}
\ArrowheadType{System Design Description}
\ArrowheadTypeShort{SysDD}
\ArrowheadVersion{5.2.2}
\ArrowheadDate{2026-10-30}
\ArrowheadAuthor{Jerker Delsing}
\ArrowheadStatus{DRAFT}
\ArrowheadContact{jerker.delsing@ltu.se}
\ArrowheadFooter{\href{www.arrowhead.eu}{www.arrowhead.eu}}
\ArrowheadSetup
%%

%% Front Page
\begin{center}
  \vspace*{1cm}
  \huge{\arrowtitle}

  \vspace*{0.2cm}
  \LARGE{\arrowtype}
  \vspace*{1cm}

  \vspace*{\fill}

  % Front Page Image

  \vspace*{1cm}
  \vspace*{\fill}

  % Front Page Abstract
  \begin{abstract}
    This document describes the design and implementation of the GraphDBManager support system of Arrowhead 5.2.2 (\texttt{ah5-GraphDBManager-java-spring}).
  \end{abstract}

  \vspace*{1cm}
\end{center}

\newpage
%%

%% Table of Contents
\tableofcontents
\newpage
%%

\section{Overview}
\label{sec:overview}

This document describes the GraphDBManager system implementation, which realizes the GraphDBManager SysD \cite{gdm-sysd}. GraphDBManager caches and persists SenML (RFC 8428) sensor and actuator data: \texttt{graphDataProxy} keeps the latest message per system and service, \texttt{graphDataHistorian} the time series, kept as RDF in Ontotext GraphDB \cite{graphdb}.
In Section \ref{Implementation}, we describe implementation details of the system.

\newpage

\section{Implementation}
\label{Implementation}

This implementation is based on the SysD document GraphDBManager \cite{gdm-sysd} and is located in \texttt{ah5-GraphDBManager-java-spring}.

\subsection{Implementation language and tools}
\begin{itemize}
\item Java 21, Spring Boot 3.4.9, Arrowhead common library 5.2.2 (private source under \texttt{vendor/}), Log4j 2; Eclipse RDF4J 5.1.2 client \cite{rdf4j} (\texttt{rdf4j-\allowbreak{}client}, HTTP repository).
\item Ontotext GraphDB Free 10.8.5 \cite{graphdb} as RDF store: one repository \texttt{graphdbmanager} (GraphDB-Free sail, \texttt{empty} ruleset, context index on). No relational database, JPA or Hibernate. Stateless request processing apart from the in-memory proxy cache.
\end{itemize}

\subsection{Independent source and build ownership}
The project folder \texttt{ah5-\allowbreak{}GraphDBManager-\allowbreak{}java-\allowbreak{}spring} is a copy of the standalone DataManager folder \cite{dm-sysd-ref} and owns its application source, Maven build, tests, executable JAR, Dockerfiles (application and GraphDB), configuration and Compose definition. Build with Java 21 and \texttt{mvn -B clean verify}; after a first build that fetches RDF4J from Maven Central, offline builds work with the populated cache.

Arrowhead common utilities and DTOs are private source copies in \texttt{vendor/}, compiled into this system by the local Maven POM. In this copy the JPA-dependent parts of the common library are removed (\texttt{common/\allowbreak{}jpa}, \texttt{Log\allowbreak{}Service}, \texttt{Log\allowbreak{}Validation}, the Log4j 2 JDBC connection factory); \texttt{Constants} holds the former entity column sizes. The removals are recorded in \texttt{vendor/\allowbreak{}PROVENANCE.json}.

The packaged artifact is \texttt{arrowhead-\allowbreak{}graphdbmanager-5.2.2.jar} (main class \texttt{eu.\allowbreak{}arrowhead.\allowbreak{}graphdbmanager.\allowbreak{}Graph\allowbreak{}DB\allowbreak{}Manager\allowbreak{}Main}).

\subsection{Functional properties implementation}

The following components (packages below \texttt{eu.\allowbreak{}arrowhead.\allowbreak{}graphdbmanager}) implement the system:

\noindent\begin{longtable}{| p{5.6cm} | p{9.6cm} |} \hline
\rowcolor{gray!33} Component & Responsibility \\ \hline \endhead
\texttt{api.\allowbreak{}http} & Controllers of graphDataProxy and graphDataHistorian. \\ \hline
\texttt{service} & Proxy cache and historian logic, SenML validation, name normalization, sysMonitor health indicator. \\ \hline
\texttt{security} & Operation resolver and authorization filter (ConsumerAuthorization \texttt{verify} per call). \\ \hline
\texttt{graphdb} & \texttt{Historian\allowbreak{}Store} interface, \texttt{Historian\allowbreak{}Pack} record and \texttt{Graph\allowbreak{}DB\allowbreak{}Historian\allowbreak{}Store} (RDF4J \texttt{HTTP\allowbreak{}Repository}). \\ \hline
\texttt{init} & Re-registration of the HTTP offers with \texttt{CERT\_\allowbreak{}AUTH}. \\ \hline
\end{longtable}

Every request passes the common filter chain: authentication filter (\texttt{authentication.\allowbreak{}policy}), optional blacklist filter, management filter (management services), the system's authorization filter, then the controller, validation and normalization, the service logic and the historian store.

\subsubsection{Data handled and stored}

\texttt{put-\allowbreak{}data} on graphDataHistorian writes one SenML pack in one RDF4J transaction into the named graph \texttt{<urn:\allowbreak{}arrowhead:\allowbreak{}graphdbmanager:\allowbreak{}historian>} of the repository. Vocabulary prefix \texttt{gdm:} = \texttt{https://\allowbreak{}arrowhead.eu/\allowbreak{}ns/\allowbreak{}graphdbmanager\#}:

\noindent\begin{longtable}{| p{5.0cm} | p{10.2cm} |} \hline
\rowcolor{gray!33} Resource / property & Content \\ \hline \endhead
\texttt{<urn:uuid:\ldots>} \texttt{a gdm:SenMLPack} & One stored pack (random UUID). \\ \hline
\texttt{gdm:systemName}, \texttt{gdm:serviceName} & Normalized endpoint (string literals). \\ \hline
\texttt{gdm:baseTime}, \texttt{gdm:minTime}, \texttt{gdm:maxTime} & Base time and time range of the pack, \texttt{xsd:double} seconds; used by the time-range query. \\ \hline
\texttt{gdm:payload} & The normalized pack as SenML JSON, returned unchanged by \texttt{get-\allowbreak{}data}. \\ \hline
\texttt{gdm:createdAt} & Storage time, \texttt{xsd:dateTime}. \\ \hline
\texttt{gdm:measurement} $\rightarrow$ \texttt{a gdm:Measurement} & One resource per named record (\texttt{<urn:uuid:\ldots\#m\textit{i}>}) with \texttt{gdm:name}, \texttt{gdm:unit} (record or base unit), absolute \texttt{gdm:time} and \texttt{gdm:value}, \texttt{gdm:stringValue} or \texttt{gdm:booleanValue}; for SPARQL users. \\ \hline
\end{longtable}

\texttt{list-\allowbreak{}systems}, \texttt{list-\allowbreak{}services} and \texttt{get-\allowbreak{}data} are SPARQL \texttt{SELECT} queries \cite{sparql11} on that graph (\texttt{DISTINCT} names; packs with \texttt{maxTime >= from} and \texttt{minTime <= to}, newest base time first, \texttt{LIMIT} count). Request values are bound to query variables as RDF literals, never concatenated into the query text. graphDataProxy keeps the last pack per endpoint in memory only.

The data models of the interfaces are Java records (of \texttt{eu.arrowhead.dto} in the common library, or of the module); the IDD documents name them without the \texttt{DTO} suffix. The records used by this system are: \texttt{SenMLRecordDTO}, \texttt{Service\allowbreak{}List\allowbreak{}Response\allowbreak{}DTO}, \texttt{System\allowbreak{}List\allowbreak{}Response\allowbreak{}DTO}.

\subsection{Non functional properties implementation}

\subsubsection{Security}

TLS and client authentication are configured with the Spring \texttt{server.\allowbreak{}ssl.\allowbreak{}*} properties (PKCS\#12 keystore and truststore, \texttt{client-\allowbreak{}auth=\allowbreak{}need}). The common library's filters implement the authentication policies; \texttt{Token\allowbreak{}Certificate\allowbreak{}Binding} binds identity tokens to the client certificate and \texttt{Crl\allowbreak{}Revocation\allowbreak{}Config} enables CRL checking. GraphDBManager authenticates to GraphDB with HTTP Basic as user \texttt{graphdbmanager}, which has only \texttt{READ\_\allowbreak{}REPO\_\allowbreak{}graphdbmanager} and \texttt{WRITE\_\allowbreak{}REPO\_\allowbreak{}graphdbmanager}.

\subsubsection{Power management}

None; the system is a server process without power states. The \texttt{sysMonitor} \texttt{power} observation always reports line supply.

\subsubsection{Internal monitoring}

Log4j 2 logs to stdout and files; there is no database log appender, so \texttt{general\allowbreak{}Management.\allowbreak{}get-\allowbreak{}log} has no stored entries. The common library's \texttt{Sys\allowbreak{}Monitor\allowbreak{}Publisher} publishes \texttt{echo}, \texttt{power} and \texttt{status} every 30 s; \texttt{echo} is \texttt{degraded} while the historian query fails.

\subsubsection{Configuration}

Configuration is read from \texttt{application.\allowbreak{}properties} and environment variables (Spring relaxed binding, e.g.\ \texttt{GRAPHDB\_\allowbreak{}URL}); secrets come from files through \texttt{SPRING\_\allowbreak{}CONFIG\_\allowbreak{}IMPORT=\allowbreak{}configtree:\allowbreak{}/\allowbreak{}run/\allowbreak{}secrets/\allowbreak{}} (\texttt{ah.\allowbreak{}graphdb-\allowbreak{}password}, \texttt{ah.\allowbreak{}store-\allowbreak{}password}, \texttt{ah.\allowbreak{}identity-\allowbreak{}password}). The properties are listed in the SysD; \texttt{general\allowbreak{}Management.\allowbreak{}get-\allowbreak{}config} returns the public ones.

\subsection{Deployment}

Two containers: the application (UID 10001, read-only root filesystem, no capabilities, HTTPS port 8467) and \texttt{graphdbmanager-\allowbreak{}graphdb}, built from \texttt{ontotext/\allowbreak{}graphdb:10.8.5} by \texttt{deployment/\allowbreak{}graphdb/\allowbreak{}Dockerfile} (UID 10001, read-only root filesystem, data volume \texttt{/\allowbreak{}opt/\allowbreak{}graphdb/\allowbreak{}home}, no host port). Its \texttt{entrypoint.sh} starts GraphDB and on every start, idempotently, creates the repository from \texttt{repository-\allowbreak{}config.ttl}, sets the admin password, creates or updates the \texttt{graphdbmanager} user, enables GraphDB security and disables free access; only then is the container healthy. Passwords are read from files and passed to curl through a temporary netrc file and standard input. The two containers share an internal network that no other system joins; GraphDB is reached over plain HTTP on it.

The local \texttt{compose.yaml} runs both on an operator-selected external cloud network (\texttt{docker compose up -d --build}). The Arrowhead-test-cloud-java test cloud is a separate deployment consumer of this primary source: there the services \texttt{graphdbmanager} and \texttt{graphdbmanager-\allowbreak{}graphdb} are added by \texttt{make up-\allowbreak{}graphdbmanager}: certificate and keystore (\texttt{generate-x509.sh}), identity (\texttt{provision-\allowbreak{}authentication-\allowbreak{}identities.sh}), AMQP access list, ConsumerAuthorization grants (\texttt{configure-\allowbreak{}graphdbmanager-\allowbreak{}policies.sh}, verifier and sysMonitor provider entries), then the two containers; the internal network \texttt{db-\allowbreak{}graphdbmanager} has a fixed subnet. ServiceRegistry, Authentication, DSO, ConsumerAuthorization and the AMQP broker remain network dependencies.

\subsection{Local modifications}
\label{sec:local}

Notes on this implementation in the workspace:
\begin{itemize}
  \item Copy of DataManager 5.2.2 (2026-10-28) with JPA/MySQL replaced by GraphDB; DataManager's \texttt{/\allowbreak{}echo} replacement by \texttt{sysMonitor} is kept.
  \item Service definitions \texttt{graphDataProxy} and \texttt{graphDataHistorian}, base path \texttt{/\allowbreak{}graphdbmanager}: ServiceRegistry's reserved service definitions bind \texttt{dataProxy} and \texttt{dataHistorian} to DataManager.
  \item \texttt{put-\allowbreak{}data} answers 204.
\end{itemize}

\subsection{Verification}

The JUnit suite covers the historian service with a mocked store, SenML validation, the proxy, the health indicator and the operation resolver and authorization filter. In the Arrowhead-test-cloud-java deployment consumer, \texttt{scripts/\allowbreak{}smoke-\allowbreak{}graphdbmanager.sh} checks, as SmokeClient and through orchestration only, for both HTTP services: \texttt{put-\allowbreak{}data}, \texttt{get-\allowbreak{}data} round trip and \texttt{list-\allowbreak{}systems}; then a SPARQL count of the stored measurement in GraphDB and the refusal (401) of anonymous GraphDB access. \texttt{smoke-test.sh} checks mTLS and identity-token access and \texttt{smoke-security.sh} the port's certificate and CRL handling. Historian data survived a restart of the GraphDB container.

\newpage

\section{Services}
\label{sec:services}

The implementation of the services is based on the SD and IDD documents referenced below.

\noindent\begin{longtable}{| p{5.4cm} | p{2.2cm} | p{2.2cm} | p{4.8cm} |} \hline
\rowcolor{gray!33}\textbf{Services produced} & \textbf{SysD ref} & \textbf{SD ref} & \textbf{IDD ref} \\ \hline \endhead
\texttt{graphDataProxy} & \cite{gdm-sysd} & \cite{gdm-graphDataProxy-sd} & \cite{gdm-graphDataProxy-http-idd} \\ \hline
\texttt{graphDataHistorian} & \cite{gdm-sysd} & \cite{gdm-graphDataHistorian-sd} & \cite{gdm-graphDataHistorian-http-idd} \\ \hline
\texttt{sysMonitor} & \cite{gdm-sysd} & \cite{gdm-sysMonitor-sd} & \cite{gdm-sysMonitor-amqp-idd} \\ \hline
&&& \\ \hline
\rowcolor{gray!33}\textbf{Services consumed} & \textbf{SysD ref} & \textbf{SD ref} & \textbf{IDD ref} \\ \hline
systemDiscovery, serviceDiscovery & \cite{core-sr-sysd} & \cite{core-sr-systemDiscovery-sd} & -- \\ \hline
identity & \cite{core-auth-sysd} & \cite{core-auth-identity-sd} & -- \\ \hline
authorization & \cite{core-ca-sysd} & \cite{core-ca-authorization-sd} & -- \\ \hline
authorization & \cite{core-ca-sysd} & \cite{core-ca-authorization-sd} & -- \\ \hline
sysMonitor (publisher) & AMQP broker (SysMonitor) & -- & -- \\ \hline
RDF4J HTTP protocol (GraphDB) & -- & -- & \cite{rdf4j,graphdb} \\ \hline
\end{longtable}


\bibliographystyle{IEEEtran}
\bibliography{bibliography}

\newpage

\section{Revision History}
\subsection{Amendments}

\noindent\begin{tabularx}{\textwidth}{| p{1cm} | p{3cm} | p{2cm} | X | p{4cm} |} \hline
\rowcolor{gray!33} No. & Date & Version & Subject of Amendments & Author \\ \hline

1 & 2026-10-28 & \arrowversion & Document created from the DataManager SysDD for the GraphDB-backed copy GraphDBManager. & Jerker Delsing \\ \hline
2 & 2026-10-30 & \arrowversion & Checked against the primary Java source tree (Eclipse-Arrowhead-java): services, operations, paths and data models unchanged. & Claude Opus 5.5 (draft) \\ \hline
\end{tabularx}

\subsection{Quality Assurance}

\noindent\begin{tabularx}{\textwidth}{| p{1cm} | p{3cm} | p{2cm} | X |} \hline
\rowcolor{gray!33} No. & Date & Version & Approved by \\ \hline

1 & & \arrowversion  &  \\ \hline

\end{tabularx}

\end{document}
